\section{3D Benchmark Results}
\label{sec:results_3d}

This section covers two experimental modes on the four canonical 3D domains.
The fixed-$k$ sweep tests all five skip factors with ten seeds per configuration.
The adaptive-$k$ results follow in Section~\ref{subsec:3d_adaptive}.

The four domains are the rectangular prism, cylinder, stacked cubes, and
L-shape prism.
The rectangular prism has uniform cooling.
The cylinder has curved surface normals.
The stacked cubes domain has an internal material interface.
The L-shape prism has a reentrant edge that concentrates the gradient.

Table~\ref{tab:3d_params} lists the physical and training parameters used
across all 3D benchmark runs.

\begin{table}[H]
\centering
\caption{Physical and training parameters for the 3D benchmark simulations.}
\label{tab:3d_params}
\small
\begin{tabular}{llc}
\toprule
Parameter & Symbol & Value \\
\midrule
Thermal conductivity        & $k_T$      & $160\,\mathrm{W/(mK)}$ \\
Volumetric heat capacity    & $\rho c_p$ & $2.4\times10^6\,\mathrm{J/(m^3K)}$ \\
Convection coefficient      & $h$        & $4000\,\mathrm{W/(m^2K)}$ \\
Initial temperature         & $T_0$      & $540\,^\circ\mathrm{C}$ \\
Water (quench) temperature  & $T_w$      & $20\,^\circ\mathrm{C}$ \\
\midrule
Adam epochs (Bayesian/TPE)  & ---        & 800 \\
Adam epochs (NSGA-II/III)   & ---        & 1\,500 \\
L-BFGS iterations per window& ---        & 100 \\
\bottomrule
\end{tabular}
\end{table}

Table~\ref{tab:summary_3d_mae} gives the ten-seed mean~$\pm$~std for all four
3D domains.
Full per-seed tables are in Appendix~\ref{app:full_results}.

\begin{table}[H]
\centering\small\setlength{\tabcolsep}{4pt}
\caption{Ten-seed mean-window MAE ($^\circ\mathrm{C}$) summary for all 3D canonical domains.
\cellcolor{RoyalBlue!20}\phantom{x} within $15\,^\circ\mathrm{C}$;
{\color{gray!60!black}grey} = above threshold.}
\label{tab:summary_3d_mae}
\begin{tabular}{@{}llccccc@{}}
\toprule
Domain & Architecture
  & \makecell{$k{=}1$\\(0\%)} & \makecell{$k{=}2$\\(50\%)} & \makecell{$k{=}3$\\(65\%)} & \makecell{$k{=}4$\\(75\%)} & \makecell{$k{=}5$\\(80\%)} \\
\midrule
\multirow{3}{*}{Rectangular 3D}
  & Bayesian/TPE & \cellcolor{RoyalBlue!20}$8.42\pm1.48$ & \cellcolor{RoyalBlue!20}$8.47\pm1.41$ & \cellcolor{RoyalBlue!20}$8.74\pm2.06$ & \cellcolor{RoyalBlue!20}$9.43\pm1.68$ & \cellcolor{RoyalBlue!20}$9.88\pm1.69$ \\
  & NSGA-II      & \cellcolor{RoyalBlue!20}$5.47\pm1.38$ & \cellcolor{RoyalBlue!20}$7.62\pm2.13$ & \cellcolor{RoyalBlue!20}$9.91\pm2.22$ & \cellcolor{RoyalBlue!20}$11.01\pm2.50$ & \cellcolor{RoyalBlue!20}$13.42\pm4.82$ \\
  & NSGA-III     & \cellcolor{RoyalBlue!20}$6.05\pm1.52$ & \cellcolor{RoyalBlue!20}$8.22\pm1.95$ & \cellcolor{RoyalBlue!20}$9.91\pm3.13$ & \cellcolor{RoyalBlue!20}$11.08\pm2.75$ & \cellcolor{RoyalBlue!20}$13.61\pm4.58$ \\
\midrule
\multirow{3}{*}{Cylinder 3D}
  & Bayesian/TPE & \cellcolor{RoyalBlue!20}$5.26\pm0.97$ & \cellcolor{RoyalBlue!20}$6.51\pm1.37$ & \cellcolor{RoyalBlue!20}$7.25\pm1.04$ & \cellcolor{RoyalBlue!20}$8.59\pm1.06$ & \cellcolor{RoyalBlue!20}$9.58\pm1.13$ \\
  & NSGA-II      & \cellcolor{RoyalBlue!20}$5.79\pm1.95$ & \cellcolor{RoyalBlue!20}$7.96\pm3.47$ & \cellcolor{RoyalBlue!20}$11.36\pm5.94$ & \cellcolor{RoyalBlue!20}$12.67\pm8.65$ & {\color{gray!60!black}$19.21\pm12.95$} \\
  & NSGA-III     & \cellcolor{RoyalBlue!20}$5.79\pm2.67$ & \cellcolor{RoyalBlue!20}$7.48\pm2.74$ & \cellcolor{RoyalBlue!20}$8.51\pm4.69$ & \cellcolor{RoyalBlue!20}$12.06\pm7.60$ & \cellcolor{RoyalBlue!20}$9.60\pm3.15$ \\
\midrule
\multirow{3}{*}{Stacked 3D}
  & Bayesian/TPE & \cellcolor{RoyalBlue!20}$10.51\pm1.40$ & \cellcolor{RoyalBlue!20}$11.35\pm2.44$ & \cellcolor{RoyalBlue!20}$11.84\pm1.72$ & \cellcolor{RoyalBlue!20}$13.03\pm2.47$ & \cellcolor{RoyalBlue!20}$14.07\pm3.13$ \\
  & NSGA-II      & \cellcolor{RoyalBlue!20}$8.69\pm3.79$ & \cellcolor{RoyalBlue!20}$11.33\pm4.15$ & \cellcolor{RoyalBlue!20}$13.31\pm5.24$ & {\color{gray!60!black}$15.24\pm4.83$} & {\color{gray!60!black}$18.19\pm5.62$} \\
  & NSGA-III     & \cellcolor{RoyalBlue!20}$8.95\pm2.71$ & \cellcolor{RoyalBlue!20}$10.96\pm4.11$ & \cellcolor{RoyalBlue!20}$12.40\pm3.44$ & \cellcolor{RoyalBlue!20}$14.29\pm4.09$ & {\color{gray!60!black}$18.81\pm6.15$} \\
\midrule
\multirow{3}{*}{L-Shape 3D}
  & Bayesian/TPE & \cellcolor{RoyalBlue!20}$10.20\pm1.48$ & \cellcolor{RoyalBlue!20}$12.10\pm2.82$ & \cellcolor{RoyalBlue!20}$12.06\pm2.57$ & \cellcolor{RoyalBlue!20}$12.19\pm2.21$ & \cellcolor{RoyalBlue!20}$13.96\pm2.26$ \\
  & NSGA-II      & \cellcolor{RoyalBlue!20}$4.54\pm1.27$ & \cellcolor{RoyalBlue!20}$7.59\pm2.87$ & \cellcolor{RoyalBlue!20}$8.83\pm2.34$ & \cellcolor{RoyalBlue!20}$10.19\pm2.15$ & {\color{gray!60!black}$15.65\pm9.32$} \\
  & NSGA-III     & \cellcolor{RoyalBlue!20}$4.41\pm1.07$ & \cellcolor{RoyalBlue!20}$7.50\pm2.07$ & \cellcolor{RoyalBlue!20}$8.94\pm3.08$ & \cellcolor{RoyalBlue!20}$10.33\pm2.04$ & \cellcolor{RoyalBlue!20}$12.26\pm2.34$ \\
\bottomrule
\end{tabular}
\end{table}

Table~\ref{tab:l2_rectangular_3d} confirms that the Rectangular Prism is the
most tractable 3D domain for all architectures. Bayesian/TPE achieves
excellent field reconstruction ($L_2 = 0.033$--$0.038$) across the entire
skip-factor sweep. It remains well below the $L_2=0.07$ excellent threshold
even at $k=5$ (80\% FEM saving). NSGA-II maintains acceptable $L_2$ through
$k=4$ but reaches $0.071$ at $k=5$, coinciding with the MAE mean exceeding the
$15\,^\circ\mathrm{C}$ threshold. NSGA-III stays below $0.07$ at all skip
factors. Seed variance is low for Bayesian/TPE ($\sigma \leq 0.008$) but
higher for NSGA-II at $k=5$ ($\sigma=0.022$). This reflects greater sensitivity
to initialisation when the prediction interval is longest.

The Cylinder results expose a pattern that repeats on the Stacked Cubes
domain. NSGA-II exhibits markedly higher inter-seed variance than both
Bayesian/TPE and NSGA-III, and its mean MAE degrades steeply with $k$.
The structural reasons for this behaviour, related to tanh saturation and the
compact architecture's loss landscape on curved and interface-rich geometries,
are analysed in Section~\ref{sec:disc_nas_transfer}.

Table~\ref{tab:l2_cylinder_3d} reveals the most pronounced architecture
divergence of all 3D domains. Bayesian/TPE maintains excellent $L_2$
reconstruction ($0.024$--$0.036$) across all skip factors, confirming that its
deeper network handles the curved surface gradients well. NSGA-II shows high
inter-seed variance at $k=1$ ($\sigma=0.020$, seed range $0.015$--$0.083$),
indicating that the compact $3{\times}153$ architecture is sensitive to
initialisation on cylindrical geometry. At $k=5$ the mean reaches $0.067$,
still within the acceptable range, but with $\sigma=0.043$ reflecting large
seed-to-seed spread. NSGA-III is more stable than NSGA-II at intermediate $k$
and stays in the excellent range through $k=4$ ($L_2=0.045$, $\sigma=0.030$).
The Cylinder is the 3D domain where the gap between Bayesian/TPE and the NSGA
variants is largest.


The Stacked Cubes domain is the most demanding 3D benchmark for NSGA-II and
NSGA-III. NSGA-II exceeds the $15\,^\circ\mathrm{C}$ threshold at $k\geq4$
(mean $15.24\,^\circ\mathrm{C}$) and NSGA-III at $k=5$ (mean
$18.81\,^\circ\mathrm{C}$). Although NSGA-II's ten-seed mean at $k=1$ is
$8.69\,^\circ\mathrm{C}$, the large inter-seed variance ($\sigma=3.79\,^\circ\mathrm{C}$)
shows that individual seeds already approach or exceed the threshold even
at the lowest skip factor. The material interface between the two cubes
creates a gradient discontinuity that is the critical difficulty. This failure
mode is not shared by the Cylinder or the L-shape. It is discussed in
Section~\ref{sec:disc_nas_transfer}.

Table~\ref{tab:l2_stacked_3d} reflects the Stacked Cubes domain as the most
challenging 3D geometry for the compact architectures. Bayesian/TPE achieves
excellent field reconstruction ($L_2 = 0.042$--$0.055$) across all skip
factors, remaining well below the $0.07$ threshold even at $k=5$. The
material interface introduces spatial complexity that raises Bayesian/TPE's
$L_2$ relative to the simpler Rectangular Prism. NSGA-II and NSGA-III both
achieve excellent $L_2$ at $k=1$ ($0.040$ and $0.037$ respectively), but with
large inter-seed spread ($\sigma \leq 0.025$). Both architectures reach the
acceptable zone at $k=4$--$5$ ($L_2 = 0.059$--$0.074$), consistent with MAE
means approaching or exceeding $15\,^\circ\mathrm{C}$ at those skip factors.
The Stacked Cubes results motivate the use of the larger Bayesian/TPE
architecture for geometries with internal material interfaces.


Table~\ref{tab:l2_lshape_3d} presents a notable reversal from the other 3D
domains. NSGA-II and NSGA-III achieve lower $L_2$ at $k=1$--$4$ than
Bayesian/TPE on the L-shape Prism. This reflects the compact architecture's
advantage on domains where the dominant thermal gradient is along a principal
axis. The L-shape prism's reentrant edge concentrates gradients in a
predictable location that the three-layer tanh networks resolve efficiently.
Bayesian/TPE's $L_2$ ($0.033$--$0.040$) nevertheless remains in the excellent
range across all $k$. At $k=5$, NSGA-II degrades sharply ($L_2=0.094$,
$\sigma=0.032$), consistent with the mean MAE exceeding $27\,^\circ\mathrm{C}$.
NSGA-III and Bayesian/TPE remain within the acceptable range at $k=5$.


Volumetric temperature field comparisons for all architectures and domains at
the lowest-error skip factor are in Appendix~\ref{app:temp_fields}
(Figure~\ref{fig:app_3d_bestk}).
Errors concentrate on curved surfaces and at material interfaces.

The per-window diagnostics reveal the same physical pattern observed in 2D.
The first prediction windows carry the highest error because the quenching
gradient is steepest at the onset of cooling. The error falls as the
temperature field approaches equilibrium. Bayesian/TPE shows the most stable
window-to-window trajectory across all four 3D domains. NSGA-II and NSGA-III
exhibit higher variance in the early windows, particularly on the Cylinder and
Stacked Cubes domains. Per-window error evolution is shown in Figure~\ref{fig:3d_l2_curve}.

% ─────────────────────────────────────────────────────────────────────────────
\subsection{Adaptive-\texorpdfstring{$k$}{k} on 3D Domains}
\label{subsec:3d_adaptive}

We applied the adaptive-$k$ controller to all four 3D domains using all three
architectures.
The controller uses a three-window burn-in at $k=1$, then promotes to a longer
window when per-window MAE stays below $\tau_\mathrm{up}=7\,^\circ\mathrm{C}$
for two consecutive windows and demotes when it exceeds
$\tau_\mathrm{down}=10\,^\circ\mathrm{C}$.

Table~\ref{tab:3d_adaptive_summary} summarises the mean-window MAE and FEM saving
for each (domain, architecture) combination.
The FEM saving is substantially lower than the 83--85\% achieved on the 2D domains
(Section~\ref{sec:results_2d_adaptive}).
The root cause is the higher base error on 3D geometries.
At $k=1$, mean per-window MAE ranges from $5$--$10\,^\circ\mathrm{C}$
(Table~\ref{tab:summary_3d_mae}), which is already close to or above the
$7\,^\circ\mathrm{C}$ promotion threshold.
The controller therefore holds $k=1$ through most of the simulation.
Bayesian/TPE is the only architecture that promotes on three of the four domains,
reaching $k=2$--$3$ and up to 35\% FEM saving on the Cylinder.
NSGA-II and NSGA-III rarely promote.
Their higher per-window error locks the controller at $k=1$ on most domains.

\begin{table}[H]
\centering\small\setlength{\tabcolsep}{4pt}
\caption{Adaptive-$k$ controller results on all four 3D domains (single seed).
Controller: $k_\mathrm{init}{=}1$, $\tau_\mathrm{up}{=}7\,^\circ\mathrm{C}$,
$\tau_\mathrm{down}{=}10\,^\circ\mathrm{C}$, 3-window burn-in.
Grey = above the $15\,^\circ\mathrm{C}$ acceptance threshold.}
\label{tab:3d_adaptive_summary}
\begin{tabular}{@{}llccc@{}}
\toprule
Domain & Architecture & MAE ($^\circ\mathrm{C}$) & FEM saving & Max $k$ \\
\midrule
\multirow{3}{*}{Rectangular}
  & Bayesian/TPE & 9.30  & 20\% & 3 \\
  & NSGA-II      & 12.06 &  5\% & 2 \\
  & NSGA-III     & 11.24 &  0\% & 1 \\
\midrule
\multirow{3}{*}{Cylinder}
  & Bayesian/TPE & 6.66           & 35\% & 3 \\
  & NSGA-II      & 13.21          &  0\% & 1 \\
  & NSGA-III     & {\color{gray!60!black}15.32} &  0\% & 1 \\
\midrule
\multirow{3}{*}{Stacked Cubes}
  & Bayesian/TPE & 8.92  & 15\% & 2 \\
  & NSGA-II      & 12.58 &  0\% & 1 \\
  & NSGA-III     & 11.23 & 10\% & 2 \\
\midrule
\multirow{3}{*}{L-Shape Prism}
  & Bayesian/TPE & 11.96 &  0\% & 1 \\
  & NSGA-II      & 13.11 & 30\% & 3 \\
  & NSGA-III     & 11.29 & 30\% & 3 \\
\bottomrule
\end{tabular}
\end{table}

For 3D domains, fixed-$k$ at $k{=}4$--$5$ (Table~\ref{tab:summary_3d_mae})
is the more practical strategy: it achieves 75--80\% FEM saving with
Bayesian/TPE while keeping mean-window MAE within the acceptance band.
The adaptive controller provides benefit on 3D only when a geometry-specific
promotion threshold is used.
This extension is described in Section~\ref{sec:conc_future}.

Figure~\ref{fig:3d_seed_bars} summarises the mean-window MAE across ten
independent seeds for all 3D benchmark configurations.  Full numerical
values are provided in the per-domain MAE tables in
Appendix~\ref{app:full_results}
(Tables~\ref{tab:cons_cylinder_3d}--\ref{tab:cons_lshape_3d}).

\begin{figure}[H]
\centering
\includegraphics[width=\textwidth]{seed_bar_3d.png}
\caption{Seed-to-seed variability in mean-window MAE for all 3D domains.
Line: ten-seed mean; band: $\pm$one standard deviation.}
\label{fig:3d_seed_bars}
\end{figure}


\begin{figure}[H]
\centering
\includegraphics[width=\textwidth]{l2_all_k_3d.png}
\caption{Per-window $L_2$ relative error over the 30-second quench for all 3D domains.
Line: ten-seed mean; band: $\pm$one standard deviation. Colour encodes $k$.}
\label{fig:3d_l2_curve}
\end{figure}
